% ======================================================================
% Smoking and Vaping Policies in Newly Constructed, State-Funded
% Multiunit Housing in California -- Methods
% Amanda Ta (San Diego State University)
%
% LaTeX edition of the author's manuscript methods section. The prose is
% unchanged from the original PDF; only the typesetting is new.
%
% To compile: pdflatex Ta_MUH_Smoking_Vaping_Policies_Methods.tex
% Requires the apa7 document class (TeX Live: tlmgr install apa7).
% ======================================================================

\documentclass[man]{apa7}

\usepackage[utf8]{inputenc}
\usepackage{hyperref}    % clickable URLs

\title{Smoking and Vaping Policies in Newly Constructed, State-Funded Multiunit Housing in California}
\shorttitle{Smoking and Vaping Policies in State-Funded MUH}
\author{Amanda Ta}
\affiliation{School of Public Health, San Diego State University}
\note{Dr. Georg E. Matt, PhD\\December 2025}

\begin{document}
\maketitle

\section{Methods}

\subsection{Study design}

This study employed a descriptive, cross-sectional research design using structured survey data to examine smoking and vaping policies in newly constructed, state-subsidized affordable multifamily housing (MUH) developments in California. The approach combined secondary analysis of administrative housing data and primary data collection through structured telephone surveys. This allows for both comprehensive coverage of the state-funded housing stock and detailed assessment of policy implementation and communication where public information was incomplete.

\subsection{Data Source and Study Population}

Project-level data were obtained from the California Tax Credit Allocation Committee (CTCAC) through the California State Treasurer's Office, which maintains a comprehensive public dataset of multifamily housing developments receiving Low-Income Housing Tax Credits (LIHTC). The dataset provides standardized, project-level information on multifamily housing developments receiving Low-Income Housing Tax Credits in California, including developments that also received California state housing tax credits. The initial dataset identified 769 newly constructed subsidized MUH developments placed in service between January 1, 2021, and December 1, 2025.

\subsection{Database Construction}

A working copy of the original CTCAC dataset was created for data collection purposes. To define the analytic sample, the dataset was filtered using a series of predetermined criteria. Only developments classified as new construction were included, and properties designated as single-room-occupancy were excluded due to their distinct housing characteristics and regulatory context. To focus specifically on housing developments receiving California taxpayer funding, properties that received only federal subsidies were excluded. After applying these criteria, the final study population consisted of $N = 281$ newly constructed MUH developments that received both state and federal housing subsidies. Additional columns were added to record property website URLs, contact phone numbers, and indicators related to smoking and vaping policy availability.

\subsection{Telephone Survey Sampling}

Primary data were collected through telephone surveys. A stratified random sampling strategy was used to select approximately 20\% of the study population ($N = 57$ properties), with strata defined by housing type (e.g., large family, senior, special needs, non-targeted housing) to ensure proportional representation across development categories. This approach was intended to enhance the generalizability of findings across different types of state-subsidized MUH developments.

\subsection{Telephone Survey Procedures}

Telephone surveys were conducted between December 4 and December 10, 2025. Calls were placed to property management staff using a standardized introduction and interview script to ensure consistency across respondents. Prior to each call, we reviewed property websites to tailor examples of relevant amenities (e.g., balconies, laundry rooms, parking garages) included in survey questions. Call attempts followed a standardized protocol, with up to three attempts made at staggered intervals before a property was classified as nonresponsive.

Survey questions assessed: (1) whether smoking or vaping was allowed inside residential units or on balconies or porches; (2) whether smoking or vaping was permitted within shared indoor spaces such as hallways, stairwells, laundry rooms, exercise facilities, or parking garages; (3) whether smoking and vaping rules were explicitly stated in lease agreements; and (4) whether any exceptions to smoking and vaping prohibitions existed. Responses were recorded using categorical response options (yes, no, or unsure) and entered into the study database.

\subsection{Analysis}

Data were analyzed using descriptive statistics to summarize smoking and vaping policy characteristics across surveyed properties. Frequencies and percentages were calculated for key outcomes, including whether smoking or vaping was permitted in residential units, balconies, and shared indoor spaces, and whether policies were included in lease agreements. Analyses were conducted for the full sample of surveyed properties ($N = 57$), with results interpreted to reflect the distribution of policies across newly constructed, state-subsidized MUH developments in California.

\end{document}
